\section{Introduction}\label{sec:intro}

How many congressional districts can a party make safe by drawing lines alone? Algorithms answer a weaker question. Given a vote map, a sampler or optimiser can \emph{find} a redistricting plan with many safe seats, but a plan that has been found says nothing about whether a better one exists: short-burst and column-generation methods are heuristics with no optimality guarantee \citep{cannon2023,gurnee2021,atlas2026}, and the exact integer programs that exist optimise compactness or county splits, not partisan outcomes \citep{validi2022,shahmizad2025}. The published dual bounds for representation objectives are for whole counties with a 10\% population tolerance \citep{fravel2026}, and coarsening the geography to make the optimisation easy \emph{restricts} the set of maps, so it yields a lower bound on what is possible, not an upper one. The consequence is that ``how many safe seats can the boundaries deliver?'' has, for real U.S.\ House geography, been answered only by lower bounds---and so has the question that follows from it, how far a given map is from what the boundaries allow.

This paper turns the question into a certification problem and answers it exactly at voting-precinct resolution for a class of realistic regimes. Our contributions are as follows.

\begin{enumerate}[leftmargin=*]
\item \textbf{A frontier and its dual.} We define the \emph{seat--safety frontier} $F(m)$---the largest number of districts that can simultaneously have a robust two-party share of at least $\tfrac12+m$---and the \emph{seat--safety spectrum} $\sigma^*(q)$, the largest achievable $q$-th best district margin. They are generalised inverses of each other (Theorem~\ref{thm:duality}), the area under $F$ is $\sum_q[\sigma^*(q)]_+$ (Corollary~\ref{cor:layer}), and both are monotone in the population tolerance, the subdivision budget and the scenario set (Proposition~\ref{prop:mono}). The single-threshold ``max seats'' objectives of earlier work are points on this curve, and at threshold $50\%$ the frontier of the two parties gives a certified \emph{seat range} for each state.
\item \textbf{Aggregated outer relaxations with a proof-carrying refinement loop.} We give a family of relaxations that aggregate precincts into cells, with exact-integer allocations, fractional-knapsack hull cuts and connectivity imposed on the quotient graph. We prove that every relaxation is valid at atomic resolution (Theorem~\ref{thm:valid}), that refinement is monotone (Theorem~\ref{thm:mono}), and that a solution splitting no cell \emph{is} an atomic plan (Theorem~\ref{thm:exact}). A counterexample-guided loop refines only the cells that a relaxed solution splits and ends with either an independently verified plan or an infeasibility proof (Theorem~\ref{thm:cegar}).
\item \textbf{A negative result and a fix.} Coarsening connectivity alone is provably (Proposition~\ref{prop:blind}) and empirically (Section~\ref{sec:negative}) unable to improve on the geography-free bound until cells are almost single precincts. A subdivision budget, counted in county splits as in the exact county-split literature, confines the looseness to a few counties at every refinement level (Proposition~\ref{prop:confine}) and makes the loop terminate quickly.
\item \textbf{Certified frontiers for real states.} For 16 states with two to five districts and both parties, using 2020 precincts, $\pm1\%$ population balance and at most $k-1$ county splits, we bracket the whole spectrum: 44 of 104 brackets are tight to half a point (median width 0.01 point in the six two-district states, 2.2--2.6 points for $3\le k\le5$), and the number of districts that can reach 50\%, 55\% and 60\% of the two-party vote is pinned exactly in 79 of 96 (state, party, threshold) cases (Section~\ref{sec:results}). Every bound is either an independently verified map or an exact-integer infeasibility proof, audited against a second solver and against brute-force enumeration; on the 20 queries that produce the upper ends of the two-district brackets the loop decides all, while exact solvers on the atomic model leave 6 (CP-SAT) or 7 (HiGHS) undecided after 300 s (Section~\ref{sec:baseline}).
\item \textbf{Seat ranges and a yardstick for enacted maps.} Across the 52 seats of these states, verified plans of the regime attain total counts of Democratic-majority districts (2020 presidential vote) from $11$ to $30$, and the proofs exclude anything outside $11$--$36$; Iowa, Nevada and New Mexico are states in which each party can draw a majority of the delegation, while in four states the delegation is the same for every plan. Rendered at precinct resolution, no enacted 118th-Congress plan reaches the certified capacity at any rank (median shortfall 5.9 points of district share), and the eight enacted plans that satisfy the regime give an independent, non-algorithmic test of the proofs, which none violates (Section~\ref{sec:enacted}).
\item \textbf{Budget, tolerance and scenarios.} Each extra county split provably raises Maine's best-district capacity by roughly a point (and turns New Hampshire Republicans from zero to one district), the population tolerance matters by at most half a point, and requiring a margin to hold across all statewide contests can cut a party's capacity by nearly ten points, so the spectrum doubles as a durability measure. The certified brackets and witness plans also serve as ground truth for scoring heuristic optimisers, which in our runs fall one to two points short in the largest cases.
\end{enumerate}

\paragraph{What this is not.} The frontier is a statement about the \emph{capacity of boundaries} given a vote geography and an explicitly stated regime. It is symmetric in the parties, is not a forecast, recommends no map, and makes no legal claim: satisfying our regime is neither necessary nor sufficient for a plan to be lawful, enacted plans were drawn under criteria our model does not include, and regimes with additional legal constraints are sub-regimes whose frontiers are, by monotonicity, bounded above by ours.
